% Realizações do projeto. Identificador estável da figura: 05.
\documentclass[aspectratio=169,11pt]{beamer}
% Estilo das figuras: mesmas fontes, tamanhos e cores da apresentação.
\usepackage[utf8]{inputenc}
\usepackage[T1]{fontenc}
\usepackage[brazilian]{babel}
\usepackage{amsmath,tikz}
\usetikzlibrary{arrows.meta,positioning,calc}
\usepackage[active,tightpage]{preview}
\PreviewEnvironment{tikzpicture}
\setlength{\PreviewBorder}{0pt}
\definecolor{ink}{HTML}{202124}
\definecolor{muted}{HTML}{666666}
\definecolor{blue}{RGB}{30,80,160}
\definecolor{green}{RGB}{30,130,60}
\definecolor{red}{RGB}{190,60,60}
\definecolor{light}{HTML}{E8E8E8}
\definecolor{amber}{HTML}{B68432}
\setbeamercolor{normal text}{fg=ink,bg=white}
\setbeamersize{text margin left=8mm,text margin right=8mm}
\setbeamertemplate{navigation symbols}{}
\tikzset{arr/.style={-{Stealth[length=2mm]},thick,muted},txt/.style={align=center,font=\normalsize},smalltxt/.style={align=center,font=\small,text=muted}}

\begin{document}
\begin{frame}
\begin{tikzpicture}
\foreach \x in {0,4.7,9.4} {
  \foreach \y in {0,-2.3} {
    \node[draw=light,rounded corners=2pt,minimum width=4.2cm,minimum height=1.55cm] at (\x,\y) {};
  }
}
\node[txt,text=blue] at (0,.25) {Revisão bibliográfica};
\node[smalltxt] at (0,-.3) {8 trabalhos};
\node[txt,text=blue] at (4.7,.25) {Modelo de planejamento};
\node[smalltxt] at (4.7,-.3) {Prazos e restrições};
\node[txt,text=blue] at (9.4,.25) {Artefato computacional};
\node[smalltxt] at (9.4,-.3) {3 estratégias + aplicação};
\node[txt,text=blue] at (0,-2.05) {Dados territoriais};
\node[smalltxt] at (0,-2.6) {GeoSaúde + amostras};
\node[txt,text=blue] at (4.7,-2.05) {Experimentos};
\node[smalltxt] at (4.7,-2.6) {Piloto + cidade + fatorial};
\node[txt,text=blue] at (9.4,-2.05) {Validação e análise};
\node[smalltxt] at (9.4,-2.6) {Versões, métricas e limites};
\end{tikzpicture}
\end{frame}
\end{document}
